%!TeX root=../persuasiontop.tex
\chapter[Chapter \thechapter]{}

\begin{figure}[t!]
\centering
\includegraphics[width=\linewidth]{nouveau13}
\end{figure}

\lettrine[lraise=0.3]{T}{he} remainder of Anne's time at Uppercross, comprehending only two days, was spent entirely at the Mansion House; and she had the satisfaction of knowing herself extremely useful there, both as an immediate companion, and as assisting in all those arrangements for the future, which, in Mr and Mrs Musgrove's distressed state of spirits, would have been difficulties.

They had an early account from Lyme the next morning. Louisa was much the same. No symptoms worse than before had appeared. Charles came a few hours afterwards, to bring a later and more particular account. He was tolerably cheerful. A speedy cure must not be hoped, but everything was going on as well as the nature of the case admitted. In speaking of the Harvilles, he seemed unable to satisfy his own sense of their kindness, especially of Mrs Harville's exertions as a nurse. <She really left nothing for Mary to do. He and Mary had been persuaded to go early to their inn last night. Mary had been hysterical again this morning. When he came away, she was going to walk out with Captain Benwick, which, he hoped, would do her good. He almost wished she had been prevailed on to come home the day before; but the truth was, that Mrs Harville left nothing for anybody to do.>

Charles was to return to Lyme the same afternoon, and his father had at first half a mind to go with him, but the ladies could not consent. It would be going only to multiply trouble to the others, and increase his own distress; and a much better scheme followed and was acted upon. A chaise was sent for from Crewkherne, and Charles conveyed back a far more useful person in the old nursery-maid of the family, one who having brought up all the children, and seen the very last, the lingering and long-petted Master Harry, sent to school after his brothers, was now living in her deserted nursery to mend stockings and dress all the blains and bruises she could get near her, and who, consequently, was only too happy in being allowed to go and help nurse dear Miss Louisa. Vague wishes of getting Sarah thither, had occurred before to Mrs Musgrove and Henrietta; but without Anne, it would hardly have been resolved on, and found practicable so soon.

They were indebted, the next day, to Charles Hayter, for all the minute knowledge of Louisa, which it was so essential to obtain every twenty-four hours. He made it his business to go to Lyme, and his account was still encouraging. The intervals of sense and consciousness were believed to be stronger. Every report agreed in Captain Wentworth's appearing fixed in Lyme.

Anne was to leave them on the morrow, an event which they all dreaded. <What should they do without her? They were wretched comforters for one another.> And so much was said in this way, that Anne thought she could not do better than impart among them the general inclination to which she was privy, and persuaded them all to go to Lyme at once. She had little difficulty; it was soon determined that they would go; go to-morrow, fix themselves at the inn, or get into lodgings, as it suited, and there remain till dear Louisa could be moved. They must be taking off some trouble from the good people she was with; they might at least relieve Mrs Harville from the care of her own children; and in short, they were so happy in the decision, that Anne was delighted with what she had done, and felt that she could not spend her last morning at Uppercross better than in assisting their preparations, and sending them off at an early hour, though her being left to the solitary range of the house was the consequence.

She was the last, excepting the little boys at the cottage, she was the very last, the only remaining one of all that had filled and animated both houses, of all that had given Uppercross its cheerful character. A few days had made a change indeed!

If Louisa recovered, it would all be well again. More than former happiness would be restored. There could not be a doubt, to her mind there was none, of what would follow her recovery. A few months hence, and the room now so deserted, occupied but by her silent, pensive self, might be filled again with all that was happy and gay, all that was glowing and bright in prosperous love, all that was most unlike Anne Elliot!

An hour's complete leisure for such reflections as these, on a dark November day, a small thick rain almost blotting out the very few objects ever to be discerned from the windows, was enough to make the sound of Lady Russell's carriage exceedingly welcome; and yet, though desirous to be gone, she could not quit the Mansion House, or look an adieu to the Cottage, with its black, dripping and comfortless veranda, or even notice through the misty glasses the last humble tenements of the village, without a saddened heart. Scenes had passed in Uppercross which made it precious. It stood the record of many sensations of pain, once severe, but now softened; and of some instances of relenting feeling, some breathings of friendship and reconciliation, which could never be looked for again, and which could never cease to be dear. She left it all behind her, all but the recollection that such things had been.

Anne had never entered Kellynch since her quitting Lady Russell's house in September. It had not been necessary, and the few occasions of its being possible for her to go to the Hall she had contrived to evade and escape from. Her first return was to resume her place in the modern and elegant apartments of the Lodge, and to gladden the eyes of its mistress.

There was some anxiety mixed with Lady Russell's joy in meeting her. She knew who had been frequenting Uppercross. But happily, either Anne was improved in plumpness and looks, or Lady Russell fancied her so; and Anne, in receiving her compliments on the occasion, had the amusement of connecting them with the silent admiration of her cousin, and of hoping that she was to be blessed with a second spring of youth and beauty.

When they came to converse, she was soon sensible of some mental change. The subjects of which her heart had been full on leaving Kellynch, and which she had felt slighted, and been compelled to smother among the Musgroves, were now become but of secondary interest. She had lately lost sight even of her father and sister and Bath. Their concerns had been sunk under those of Uppercross; and when Lady Russell reverted to their former hopes and fears, and spoke her satisfaction in the house in Camden Place, which had been taken, and her regret that Mrs Clay should still be with them, Anne would have been ashamed to have it known how much more she was thinking of Lyme and Louisa Musgrove, and all her acquaintance there; how much more interesting to her was the home and the friendship of the Harvilles and Captain Benwick, than her own father's house in Camden Place, or her own sister's intimacy with Mrs Clay. She was actually forced to exert herself to meet Lady Russell with anything like the appearance of equal solicitude, on topics which had by nature the first claim on her.

There was a little awkwardness at first in their discourse on another subject. They must speak of the accident at Lyme. Lady Russell had not been arrived five minutes the day before, when a full account of the whole had burst on her; but still it must be talked of, she must make enquiries, she must regret the imprudence, lament the result, and Captain Wentworth's name must be mentioned by both. Anne was conscious of not doing it so well as Lady Russell. She could not speak the name, and look straight forward to Lady Russell's eye, till she had adopted the expedient of telling her briefly what she thought of the attachment between him and Louisa. When this was told, his name distressed her no longer.

Lady Russell had only to listen composedly, and wish them happy, but internally her heart revelled in angry pleasure, in pleased contempt, that the man who at twenty-three had seemed to understand somewhat of the value of an Anne Elliot, should, eight years afterwards, be charmed by a Louisa Musgrove.

The first three or four days passed most quietly, with no circumstance to mark them excepting the receipt of a note or two from Lyme, which found their way to Anne, she could not tell how, and brought a rather improving account of Louisa. At the end of that period, Lady Russell's politeness could repose no longer, and the fainter self-threatenings of the past became in a decided tone, <I must call on Mrs Croft; I really must call upon her soon. Anne, have you courage to go with me, and pay a visit in that house? It will be some trial to us both.>

Anne did not shrink from it; on the contrary, she truly felt as she said, in observing—

<I think you are very likely to suffer the most of the two; your feelings are less reconciled to the change than mine. By remaining in the neighbourhood, I am become inured to it.>

She could have said more on the subject; for she had in fact so high an opinion of the Crofts, and considered her father so very fortunate in his tenants, felt the parish to be so sure of a good example, and the poor of the best attention and relief, that however sorry and ashamed for the necessity of the removal, she could not but in conscience feel that they were gone who deserved not to stay, and that Kellynch Hall had passed into better hands than its owners'. These convictions must unquestionably have their own pain, and severe was its kind; but they precluded that pain which Lady Russell would suffer in entering the house again, and returning through the well-known apartments.

In such moments Anne had no power of saying to herself, <These rooms ought to belong only to us. Oh, how fallen in their destination! How unworthily occupied! An ancient family to be so driven away! Strangers filling their place!> No, except when she thought of her mother, and remembered where she had been used to sit and preside, she had no sigh of that description to heave.

Mrs Croft always met her with a kindness which gave her the pleasure of fancying herself a favourite, and on the present occasion, receiving her in that house, there was particular attention.

The sad accident at Lyme was soon the prevailing topic, and on comparing their latest accounts of the invalid, it appeared that each lady dated her intelligence from the same hour of yestermorn; that Captain Wentworth had been in Kellynch yesterday (the first time since the accident), had brought Anne the last note, which she had not been able to trace the exact steps of; had staid a few hours and then returned again to Lyme, and without any present intention of quitting it any more. He had enquired after her, she found, particularly; had expressed his hope of Miss Elliot's not being the worse for her exertions, and had spoken of those exertions as great. This was handsome, and gave her more pleasure than almost anything else could have done.

As to the sad catastrophe itself, it could be canvassed only in one style by a couple of steady, sensible women, whose judgements had to work on ascertained events; and it was perfectly decided that it had been the consequence of much thoughtlessness and much imprudence; that its effects were most alarming, and that it was frightful to think, how long Miss Musgrove's recovery might yet be doubtful, and how liable she would still remain to suffer from the concussion hereafter! The Admiral wound it up summarily by exclaiming—

<Ay, a very bad business indeed. A new sort of way this, for a young fellow to be making love, by breaking his mistress's head, is not it, Miss Elliot? This is breaking a head and giving a plaster, truly!>

Admiral Croft's manners were not quite of the tone to suit Lady Russell, but they delighted Anne. His goodness of heart and simplicity of character were irresistible.

<Now, this must be very bad for you,> said he, suddenly rousing from a little reverie, <to be coming and finding us here. I had not recollected it before, I declare, but it must be very bad. But now, do not stand upon ceremony. Get up and go over all the rooms in the house if you like it.>

<Another time, Sir, I thank you, not now.>

<Well, whenever it suits you. You can slip in from the shrubbery at any time; and there you will find we keep our umbrellas hanging up by that door. A good place is not it? But,> (checking himself), <you will not think it a good place, for yours were always kept in the butler's room. Ay, so it always is, I believe. One man's ways may be as good as another's, but we all like our own best. And so you must judge for yourself, whether it would be better for you to go about the house or not.>

Anne, finding she might decline it, did so, very gratefully.

<We have made very few changes either,> continued the Admiral, after thinking a moment. <Very few. We told you about the laundry-door, at Uppercross. That has been a very great improvement. The wonder was, how any family upon earth could bear with the inconvenience of its opening as it did, so long! You will tell Sir Walter what we have done, and that Mr Shepherd thinks it the greatest improvement the house ever had. Indeed, I must do ourselves the justice to say, that the few alterations we have made have been all very much for the better. My wife should have the credit of them, however. I have done very little besides sending away some of the large looking-glasses from my dressing-room, which was your father's. A very good man, and very much the gentleman I am sure: but I should think, Miss Elliot,> (looking with serious reflection), <I should think he must be rather a dressy man for his time of life. Such a number of looking-glasses! oh Lord! there was no getting away from one's self. So I got Sophy to lend me a hand, and we soon shifted their quarters; and now I am quite snug, with my little shaving glass in one corner, and another great thing that I never go near.>

Anne, amused in spite of herself, was rather distressed for an answer, and the Admiral, fearing he might not have been civil enough, took up the subject again, to say—

<The next time you write to your good father, Miss Elliot, pray give him my compliments and Mrs Croft's, and say that we are settled here quite to our liking, and have no fault at all to find with the place. The breakfast-room chimney smokes a little, I grant you, but it is only when the wind is due north and blows hard, which may not happen three times a winter. And take it altogether, now that we have been into most of the houses hereabouts and can judge, there is not one that we like better than this. Pray say so, with my compliments. He will be glad to hear it.>

Lady Russell and Mrs Croft were very well pleased with each other: but the acquaintance which this visit began was fated not to proceed far at present; for when it was returned, the Crofts announced themselves to be going away for a few weeks, to visit their connexions in the north of the county, and probably might not be at home again before Lady Russell would be removing to Bath.

So ended all danger to Anne of meeting Captain Wentworth at Kellynch Hall, or of seeing him in company with her friend. Everything was safe enough, and she smiled over the many anxious feelings she had wasted on the subject.